% Лекция 17. Источник: Тервер(Лекция)_202611112317_42212.pdf, с. 1.
\clearpage
\begingroup
\tcbset{myboxstyle/.append style={breakable=false}}
\everymath{\textstyle}
\subsection{Разложение функции распределения}

\begin{definition}[Сингулярное непрерывное распределение]
    Функция распределения $F$ называется \textbf{сингулярной непрерывной},
    если она непрерывна и
    \[
        F'(x)=0
        \quad\text{почти всюду относительно меры Лебега}.
    \]
\end{definition}

\begin{theorem}[Разложение функции распределения]
    Любую функцию распределения можно представить в виде
    \[
        F(x)=\alpha_1F_1(x)+\alpha_2F_2(x)+\alpha_3F_3(x),
    \]
    где $F_1$ --- дискретная, $F_2$ --- абсолютно непрерывная,
    $F_3$ --- сингулярная непрерывная функция распределения, причём
    \[
        \alpha_1,\alpha_2,\alpha_3\geq0,
        \qquad \alpha_1+\alpha_2+\alpha_3=1.
    \]
\end{theorem}
\endgroup
